\documentclass{report}
\usepackage{amsmath,amssymb}
\setlength{\parindent}{0mm}
\setlength{\parskip}{1em}
\begin{document}
\begin{center}
\rule{6in}{1pt} \
{\large Osni Marques \\
{\bf Tuning the Coarse Space Construction in a Spectral AMG Solver}}

Lawrence Berkeley National Laboratory \\ 1 Cyclotron Road \\ MS 50F-1650 \\ Berkeley \\ CA 94720-8139 \\ USA
\\
{\tt oamarques@lbl.gov}\\
Alex  Druinsky\\
Sherry  Li\\
Andrew Barker\\
Delyan Kalchev\\
Panayot Vassilevski\\
	
	     \end{center}

In this presentation we will discuss strategies for computing subsets of
eigenvectors of matrices corresponding to subdomains
of finite element meshes achieving compromise between two contradicting
goals. The subset of eigenvectors is required in the construction of
coarse spaces used in algebraic multigrid methods (AMG) as well as in
certain domain decomposition (DD) methods. The quality of the coarse
spaces depends on the number of eigenvectors, which improves the
approximation properties of the coarse space and the latter impacts the
overall performance and convergence of the associated AMG or DD
algorithms. However, a large number of eigenvectors reflects negatively
the sparsity of the corresponding coarse matrices, which can become
fairly dense. The sparsity of the coarse matrices can be controlled to a
certain extent by the size of the subdomains (union of finite elements)
referred to as agglomerates. However, if the size of the agglomerates is
too large then the cost of the eigensolvers increases and eventually can
become unacceptable for the purpose of constructing the AMG or DD
solvers. This presentation discusses strategies to optimize the solution
of the partial eigenproblems of interest. In particular, we examine
direct and iterative eigensolvers for computing those subsets. Our
experiments with synthetic meshes and with a well-known model of an
oil-reservoir simulation benchmark indicate that iterative eigensolvers
can lead to significant improvement in the overall performance of an AMG
solver that exploits such spectral construction of coarse spaces.


\end{document}
